\documentclass[10pt,twocolumn,letterpaper]{article}

\usepackage{cvpr}
\usepackage{times}
\usepackage{epsfig}
\usepackage{graphicx}
\usepackage{amsmath}
\usepackage{amssymb}
\usepackage{subcaption}
\usepackage{listings}
\usepackage{url}
\usepackage{verbatim}

% Include other packages here, before hyperref.

% If you comment hyperref and then uncomment it, you should delete
% egpaper.aux before re-running latex.  (Or just hit 'q' on the first latex
% run, let it finish, and you should be clear).
%\usepackage[pagebackref=true,breaklinks=true,letterpaper=true,colorlinks,bookmarks=false]{hyperref}

\cvprfinalcopy % *** Uncomment this line for the final submission

\def\cvprPaperID{****} % *** Enter the CVPR Paper ID here
\def\httilde{\mbox{\tt\raisebox{-.5ex}{\symbol{126}}}}

% Pages are numbered in submission mode, and unnumbered in camera-ready
\ifcvprfinal\pagestyle{empty}\fi
\begin{document}

%%%%%%%%% TITLE
\title{Parallel Image Stitching}

\author{Davide Lombardi\\
E-mail address\\
{\tt\small davide.lombardi2@edu.unifi.it}
% For a paper whose authors are all at the same institution,
% omit the following lines up until the closing ``}''.
% Additional authors and addresses can be added with ``\and'',
% just like the second author.
% To save space, use either the email address or home page, not both
}

\maketitle
\thispagestyle{empty}

%%%%%%%%% ABSTRACT
\begin{abstract}
Image stitching is a computer vision task that combines multiple 
images with overlapping fields of view to produce a high-resolution panorama. 
However, the standard sequential pipeline (that comprises feature extraction, matching, homography estimation, warping, and re-extraction) 
is computationally expensive and scales poorly with large image sets. 
This project explores the parallelization of the image stitching pipeline in Python, 
evaluating multiple concurrency paradigms to overcome language-specific bottlenecks such as 
the Global Interpreter Lock (GIL) and Inter-Process Communication (IPC) serialization overhead. 
Several parallel architectures, including Task Parallelism and Data Parallelism, 
were implemented and benchmarked against a sequential baseline. 
Experimental results indicate that while isolated phases 
 like SIFT extraction scale efficiently (achieving up to a 2.8x speedup using process pools), 
 the overall aggregate performance is heavily constrained by strict sequential dependencies, 
 primarily the feature re-extraction phase on the continuously updating panorama. 
\end{abstract}



%%%%%%%%% BODY TEXT
\noindent\large\textbf{Future Distribution Permission}\\
\indent The author(s) of this report give permission for this document to be distributed to Unifi-affiliated students taking future courses.

\section{Introduction}
Image stitching is an important process in computer vision, widely utilized to construct 
high-resolution panoramic images from a sequence of overlapping photographs. 
This technique finds extensive application in fields such as geospatial mapping, 
autonomous driving, virtual reality, and consumer photography. 
Despite its widespread adoption, the creation of a seamless panorama from multiple independent 
images demands substantial computational resources. Each stage of the pipeline introduces 
unique mathematical and algorithmic challenges that scale poorly when processed sequentially, 
particularly as image resolutions and dataset sizes increase.

The main objective of this work is to analyze the performance bottlenecks inherent to 
the standard sequential image stitching pipeline and evaluate the efficacy of various parallelization strategies in Python. 
This section delineates the theoretical foundations and functional role of 
each phase within the stitching framework.

\subsection{Feature Extraction}
The initial phase of the pipeline involves identifying distinctive points within each input image
 that remain invariant to scale, rotation, and illumination changes. In this framework, 
 the Scale-Invariant Feature Transform (SIFT) algorithm is employed. SIFT detects local extrema 
 in scale-space by computing the Difference of Gaussians (DoG), defined mathematically as 
 $D(x, y, \sigma) = L(x, y, k\sigma) - L(x, y, \sigma)$, where $L$ is the image convolved with a Gaussian blur. 
For each keypoint, SIFT then computes a 128-dimensional descriptor. 
 Because feature extraction operates independently on each discrete image file, 
 this phase represents an ideal candidate for data-level parallelization across multiple worker processes.

\subsection{Feature Matching}
Once keypoints and descriptors are computed for the input images, correspondence must be 
established between adjacent frames. Feature matching computes the distance 
(typically using the L2 norm for SIFT) between descriptors in different images to identify overlapping regions.
 This implementation utilizes a Fast Library for Approximate Nearest Neighbors (FLANN) matcher combined with Lowe's ratio
 test to filter out ambiguous and spurious matches. 
 Formally, a descriptor $d_A$ is matched to its nearest neighbor $d_{B1}$ only if the ratio of distances to the first and second nearest neighbors ($d_{B2}$) satisfies 
 $\frac{||d_A - d_{B1}||_2}{||d_A - d_{B2}||_2} < \tau$, with $\tau$ typically set around $0.7$.

\subsection{Homography Estimation} 
The next step involves aligning the matched keypoints from the source image to the target image.
This is achieved by estimating a $3\times3$ homography matrix ($H$) 
using the Random Sample Consensus (RANSAC) algorithm. 
Mathematically, $H$ projects a keypoint $p = [x, y, 1]^T$ from the source image to the target coordinate space $p' = [x', y', 1]^T$ 
via the linear equation $p' \sim H p$, where the equality is up to a non-zero scale factor.
RANSAC iteratively selects a minimal subset of matched points to estimate 
the transformation model, separating valid inliers from outlier matches caused 
by false correspondences. This phase is characterized by an iterative, 
non-deterministic execution time that depends heavily on the quality of the initial feature matches.

\subsection{Warping and Blending}
After determining the homography matrices, the target images are transformed (warped) into the spatial 
domain of the reference frame or the evolving canvas. 
Warping requires computing backward coordinate mapping and pixel interpolation. 
To prevent visible seams, exposure mismatches, and ghosting artifacts in the overlapping regions, 
a blending mask or multi-band blending algorithm is applied. 
In a standard linear alpha blending, the final pixel intensity is computed as 
$I_{blend}(x,y) = \alpha(x,y) I_{ref}(x,y) + (1 - \alpha(x,y)) I_{warped}(x,y)$, where $\alpha \in [0,1]$ dictates the pixel contribution.
This phase involves heavy matrix operations and pixel-level manipulation, 
making memory bandwidth and array serialization significant factors during parallel execution.

\subsection{Feature Re-extraction (Re-ext)}
In a linear or progressive stitching pipeline, a severe structural bottleneck arises 
when appending subsequent frames. Once an image is warped and blended onto the main canvas, 
the existing keypoint coordinate system is altered, and new overlapping regions are created. 
Then, features must be re-extracted from the updated panorama to facilitate matching with the next incoming frame. 
The computational cost of this phase escalates dramatically, because the dimensions of the canvas grow progressively with each integrated image.
 If the initial image area is $A_0$, the panorama area after $n$ iterations expands to $A_n \approx A_0 + \sum_{i=1}^n \Delta A_i$,
 causing the time complexity of re-extraction to monotonically increase over time.
 Moreover, strict data dependencies require the completion of the current blending step before the next re-extraction can begin, imposing a rigid sequential constraint on the architecture.

\subsection{Pipeline Execution Overview}
To illustrate the complete workflow, the framework processes a sequential dataset of multiple overlapping images and aligns them into a single coherent panorama. The structural transformation from the independent inputs to the final synthesized output is shown in Figure~\ref{fig:stitching_overview}.

\begin{comment}

\begin{figure}[htbp]
    \centering
    % PRIMA RIGA DI INPUT: Image 1 & Image 2 affiancate nella colonna
    \begin{subfigure}[b]{0.48\linewidth}
        \centering
        \includegraphics[width=\linewidth]{images/0000.jpg}
        \caption{Image 1}
    \end{subfigure}
    \hfill
    \begin{subfigure}[b]{0.48\linewidth}
        \centering
        \includegraphics[width=\linewidth]{images/0001.jpg}
        \caption{Image 2}
    \end{subfigure}
    
    \vspace{0.15cm} % Piccolo spazio verticale tra le due righe di input
    
    % SECONDA RIGA DI INPUT: Image 3 & Image 4 affiancate nella colonna
    \begin{subfigure}[b]{0.48\linewidth}
        \centering
        \includegraphics[width=\linewidth]{images/0002.jpg}
        \caption{Image 3}
    \end{subfigure}
    \hfill
    \begin{subfigure}[b]{0.48\linewidth}
        \centering
        \includegraphics[width=\linewidth]{images/0003.jpg}
        \caption{Image 4}
    \end{subfigure}
    
    \vspace{0.3cm} % Spazio verticale prima del panorama
    
    % RIGA INFERIORE: Il panorama finale a larghezza colonna piena
    \begin{subfigure}[b]{1.0\linewidth}
        \centering
        \includegraphics[width=\linewidth]{images/final_panorama.jpg}
        \caption{Final stitched panoramic output}
    \end{subfigure}
    
\caption{Overview of the image stitching process: a sequence of independent overlapping input frames (top, shown with a sample of four images) are progressively aligned, warped, and blended to produce the unified final panorama (bottom).}    \label{fig:stitching_overview}
\end{figure}

\end{comment}

\section{Implementation Details}
This section details the experimental environment, 
the preprocessing strategy applied to the drone imagery dataset, 
and the specific concurrent architectures designed to accelerate the image stitching pipeline in Python.

\subsection{Dataset and Preprocessing Setup}

For the experimental validation of the framework, the OpenDroneMap dataset \cite{odmAukermanDataset} is used, 
which contains high-resolution aerial imagery (4896 × 3672 pixels) captured via UAVs. 
All images are downsampled to 50\% of their original resolution: this choice is strictly dictated by the necessity to compress the memory footprint and bound the initial SIFT computational complexity.

Progressive image stitching faces a major issue when processing frames in strict chronological order. 
If two consecutive images fail to align (usually due to low-texture areas or sharp changes in perspective), the main canvas stops updating. 
As a result, the system prematurely drops all remaining images, even if they share valid overlaps. 
To resolve this weakness without changing the core stitching algorithm, 
a separate diagnostic and optimization framework using \texttt{window\_diagnostics.py} and \texttt{reorder\_windows.py} is employed.

Operating within a localized sliding execution window, 
this preprocessing routine executes a brute-force pairwise analysis. 
Features are extracted once per image to compute the RANSAC inlier count for all possible $C(n,2)$ combinations, 
mapping the visual connectivity matrix of the local subset. 
Through this heuristic analysis, the algorithm identifies the "hub" image, defined as the frame exhibiting the highest degree of valid correspondences and total inlier density.
 Consequently, files are physically reordered and renamed on disk. 
 This layout ensures that each sliding window initiates its progressive blending sequence from its optimal structural
  anchor, preventing recoverable topological data from being lost due to localized tracking loss.

\subsection{Parallelization Scope and System Constraints}
The design of an effective concurrent architecture in Python is strictly governed by loop-carried dependencies and the rigid execution bottlenecks of the CPython interpreter. 
Driven by the fundamentally distinct computational demands of each processing stage, 
the achievable parallelization scope shifts considerably throughout the pipeline:

\begin{itemize}
   \item \textbf{Parallelizable phases:} loading images is an I/O-bound task,
    making it a perfect fit for multithreading. Since Python releases the GIL
    during file operations, multiple threads can read from disk at the same
    time without blocking each other. By contrast, SIFT feature extraction is
    a heavy CPU-bound task: since features are computed independently
    for each image, this step is highly parallel. A pool of worker \emph{processes} allows the system to completely bypass
    the GIL instead of just dealing with it. Finally, the warping and
    blending stage is parallelized differently: rather than splitting the
    workload per image, the output canvas is divided into horizontal tiles.
    These tiles are blended concurrently, exploiting data parallelism while
    merging two images.

    \item \textbf{Non-parallelizable phases (sequential bottleneck):} feature
    matching, homography estimation (RANSAC), and canvas blending form a
    strictly sequential chain. Each stitch operates on the panorama produced by
    the previous one, so stitch $i$ cannot begin before stitch $i-1$ has
    committed its result to the canvas. Feature re-extraction inherits this
    constraint, since it must run on the newly merged canvas rather than on the
    original input images. Given that the canvas grows monotonically with each merge
    , this re-extraction step becomes an
    increasingly expensive, strictly serial bottleneck that cannot be
    parallelized within a simple left-to-right (linear-fold) alignment
    strategy, motivating the alternative merge strategies discussed later
    (e.g.\ the tree-based MapReduce fold).
\end{itemize}

To systematically evaluate the trade-offs between concurrency paradigms,
language-level overheads, and synchronization costs, six distinct architectural
configurations were implemented and benchmarked against each other.

\subsection{Sequential Baseline (\texttt{sequential.py})}
This configuration serves as the experimental control: the entire stitching
pipeline (image loading, SIFT extraction, matching, homography estimation,
warping, blending, and re-extraction) runs on a single thread within a single
process, with no concurrency primitives of any kind. Its purpose is not
performance but \emph{measurement}: by isolating the raw, unparallelized cost of
each phase, it establishes the reference timings against which every other
architecture's speedup is computed.

Concretely, this script processes the workload in two distinct stages: it 
first extracts SIFT features for all input images upfront, and then initiates 
the iterative left-to-right blending loop. This implementation serves as a 
direct practical manifestation of the progressive re-extraction bottleneck 
described previously.

\begin{lstlisting}[language=Python, basicstyle=\small\ttfamily, frame=single, breaklines=true, breakatwhitespace=true, caption={Core sequential fold in \texttt{sequential.py}: each iteration strictly depends on the panorama state produced by the previous one.}]
base_image = images[0]
base_kp, base_des = kp_list[0], des_list[0]

for i in range(1, len(images)):
    matches = match_features(base_des, des_list[i])
    if len(matches) < 4:
        continue

    H = estimate_homography(base_kp, kp_list[i], matches)
    if H is None:
        continue

    base_image = warp_and_blend(base_image, images[i], H)

    gray = cv2.cvtColor(base_image, cv2.COLOR_BGR2GRAY)
    base_kp, base_des = sift.detectAndCompute(gray, None)
\end{lstlisting}

Every other pipeline evaluated in this study preserves this same
match~$\rightarrow$~homography~$\rightarrow$~warp~$\rightarrow$~re-extract
dependency chain unchanged; what differs between architectures is exclusively
\emph{how} the independent, non-sequential work around this chain is
scheduled onto threads or processes.

\subsection{Parallel Pipeline (\texttt{parallel.py})}
This configuration uses a targeted approach to concurrency. Instead of 
applying a single parallel method to the entire pipeline, it addresses 
three specific bottlenecks independently. The choice between using threads 
or processes is decided case by case, depending on whether the task is 
I/O-bound or CPU-bound, and whether the GIL restricts its performance.

\paragraph{Threaded I/O loading.} Image loading is handled by a
\texttt{ThreadPoolExecutor}. Since \texttt{cv2.imread} mostly waits for
the disk, Python releases the GIL during each read operation. This allows
multiple threads to load files at the same time without blocking each
other. Multiple processes are not needed here because reading files is
an I/O-bound task, not a heavy CPU computation.

\paragraph{Process-pool SIFT extraction.} SIFT feature extraction is a
heavy CPU-bound task. Using threads here would not help because Python's
GIL prevents true parallel execution across multiple images, then a
\texttt{ProcessPoolExecutor} is used. This gives each worker its own
Python interpreter, bypassing the GIL completely. However, this creates
a data transfer issue: \texttt{multiprocessing} uses "pickling" to
serialize and move data between processes, but OpenCV
\texttt{cv2.KeyPoint} objects cannot be pickled natively. To fix this,
each worker breaks down the keypoints into simple Python tuples before
returning them, and the main process rebuilds them into
\texttt{cv2.KeyPoint} objects later. Finally, since spawning new OS
processes is slow, the process pool is created only once at the
beginning and reused for all windows.

\paragraph{Tiled blending.} The final stage uses a third type of
parallelism. Instead of parallelizing across images or I/O requests,
the blending step parallelizes within a single merge operation.
Because alpha-blending relies on NumPy array arithmetic, Python
releases the GIL during these vectorized operations. Therefore, a
\texttt{ThreadPoolExecutor} is highly efficient and avoids the high
overhead of processes, even though the computation is numerically
heavy. The canvas is divided into horizontal tiles based on the
number of available CPU cores. Each tile is blended independently
by a thread, and the final results are reassembled vertically using
\texttt{np.vstack}. Additionally, this stage protects against bad
homography matrices. Before allocating memory, the expected canvas
size is checked against a set limit of the combined input area. If
the size is too large, an error is raised early to prevent
out-of-memory crashes caused by an unstable RANSAC result.

Crucially, these optimizations do not change the core sequential
dependency chain (matching, homography, warping, and re-extraction).
This chain remains strictly serial because tiling only parallelizes a
single merge operation internally; it does not allow multiple separate
merges to run at the same time.

\begin{lstlisting}[language=Python, basicstyle=\small\ttfamily, frame=single, breaklines=true, breakatwhitespace=true, caption={Process-boundary serialization workaround for \texttt{cv2.KeyPoint} in \texttt{parallel.py}: each worker returns primitive tuples, which the parent process reconstructs into native OpenCV objects.}]
def extract_single_image_features(img):
    cv2.setNumThreads(1)  # avoid oversubscription across worker processes
    sift = cv2.SIFT_create(nfeatures=8000)
    gray = cv2.cvtColor(img, cv2.COLOR_BGR2GRAY)
    kp, des = sift.detectAndCompute(gray, None)

    # cv2.KeyPoint is a C++ binding and cannot be pickled directly;
    # decompose it into a tuple of primitives for cross-process transfer.
    kp_serialized = [
        (p.pt, p.size, p.angle, p.response, p.octave, p.class_id)
        for p in kp
    ]
    return kp_serialized, des

# Back in the parent process
results = list(process_executor.map(extract_single_image_features, images))
keypoints_list = [
    [cv2.KeyPoint(pt[0], pt[1], size, angle, response, octave, class_id)
     for pt, size, angle, response, octave, class_id in kp_serialized]
    for kp_serialized, _ in results
]
\end{lstlisting}

\subsection{Producer-Consumer Paradigm}

This variant decouples the two structurally different workloads present in
the stitching pipeline: SIFT feature extraction, which has no dependency on
the state of the growing panorama, and the sequential fold chain
(matching \(\rightarrow\) homography estimation \(\rightarrow\) warping
\(\rightarrow\) re-extraction), whose steps are strictly ordered because
each stitch operates on the panorama produced by the previous one. Rather
than executing these two workloads back-to-back, the producer-consumer
pattern overlaps them in time, so that feature extraction for image
\(i+1\) proceeds concurrently with the fold operations being applied to
image \(i\).

\paragraph{Architecture.}
A single \emph{producer} thread submits every image's feature-extraction
job to a warm \texttt{ProcessPoolExecutor} at the start of the window,
allowing the pool to schedule extraction across all available cores. As
each future resolves, the producer deserialises the returned keypoints and
pushes the tuple \((\text{index}, \text{image}, \text{keypoints},
\text{descriptors})\) onto a bounded \texttt{queue.Queue}, strictly in
image order. A \emph{consumer}, running on the main thread, blocks on this
queue and, as soon as an item becomes available, drives the sequential
fold.

It is worth underlining that producer and consumer live in the \emph{same}
operating-system process, communicating through a standard
\texttt{queue.Queue} rather than a \texttt{multiprocessing.Queue}. Because
both ends share the same address space, handing an item from producer to
consumer only requires passing an object reference, with no serialisation
cost. Serialisation (pickling) is incurred exactly once per image, only at
the process boundary between the producer thread and the
\texttt{ProcessPoolExecutor} worker that performs the actual SIFT
extraction (and on the way back, for the small keypoint/descriptor
payload).

\paragraph{Backpressure and termination.} The queue has a fixed maximum 
size (\texttt{queue\_depth}, set to two items by default), which causes 
the \texttt{queue.put()} method to block automatically if the producer 
runs faster than the consumer. This mechanism slows down the producer 
and limits the number of fully-extracted images held in memory at any 
time. Once all images are sent, the producer places a special "sentinel" 
value into the queue, allowing the consumer to recognize it and exit 
its loop before the producer thread is joined and the executors are shut 
down.

\paragraph{Combining task and data parallelism.}
This variant layers a second, orthogonal form of parallelism on top of the
producer-consumer overlap. The warping and blending step invokes
\texttt{warp\_and\_blend\_tiling}, which partitions the warped canvas into
horizontal strips and blends them concurrently on a
\texttt{ThreadPoolExecutor}. The producer-consumer overlap is
\emph{task} parallelism  (different pipeline stages executing at the
same time) while the tiled blend is \emph{data} parallelism (a
single operation split across a partition of its input). Because both the
process pool (extraction) and the thread pool (blending) can be active
simultaneously, they compete for the same physical cores.

\begin{lstlisting}[language=Python, basicstyle=\footnotesize\ttfamily, breaklines=true, frame=single, breakatwhitespace=true, xrightmargin=1em, caption={Producer thread: extraction dispatch and handoff via a bounded queue.}, label={lst:producer-consumer}]
def _producer(images, process_executor, out_queue):
    # Dispatch every extraction job to the warm process pool at once
    futures = [process_executor.submit(_extract_worker, img)
               for img in images]

    for i, fut in enumerate(futures):
        # Blocks until image i's extraction result is ready, in order
        kp_serialized, des = fut.result()
        kp = _deserialize_kp(kp_serialized)

        # Hands off a reference, not a copy: producer and consumer
        # share the same process. Blocks if the queue is full
        # (backpressure).
        out_queue.put((i, images[i], kp, des))

    # Signal completion to the consumer
    out_queue.put(_SENTINEL)
\end{lstlisting}

\subsection{MapReduce Paradigm}

Functional decomposition principles inspire this fourth architectural variant. Unlike the
sequential, parallel, and producer-consumer pipelines discussed above,
 the MapReduce variant restructures the
computation as a binary merge tree, directly targeting the bottleneck identified by
Amdahl's law: in a linear fold, stitching image $i$ can only begin once image $i-1$ has been
merged, so the fold itself never parallelizes regardless of the number of available cores.

\paragraph{Map phase.}
The input dataset is first partitioned into $n$ independent image indices, each dispatched to
a worker process in a \texttt{ProcessPoolExecutor}. Because SIFT keypoint extraction operates
on a single raw image with no dependency on any panorama state, this phase is
\emph{embarrassingly parallel}: all $n$ extractions proceed concurrently, bounded only by the
number of physical cores available. Each worker returns a self-contained \emph{node}, consisting of
the original image together with its serialized keypoints and descriptor matrix, ready to
enter the reduction phase.

\paragraph{Reduce phase.}
The subsequent phase departs from linear accumulation entirely. Instead of a single running
panorama, nodes are combined pairwise in a hierarchical reduction tree:

\begin{itemize}
    \item \textbf{Level 1}: adjacent pairs $(0,1), (2,3), (4,5), \ldots$ are stitched
    concurrently, producing $\lceil n/2 \rceil$ partial panoramas.
    \item \textbf{Level 2}: pairs of partial panoramas from level 1 are stitched concurrently,
    producing $\lceil n/4 \rceil$ nodes.
    \item[] \ldots the process repeats until a single root node (the final panorama)
    remains.
\end{itemize}

Each processing level only waits for the level directly 
below it. Within the same level, however, all image pairs are 
completely independent and can be processed in parallel by different 
worker processes. This creates a tree structure with only 
$\lceil \log_2 n \rceil$ steps, which is much faster than the $n-1$ 
steps required by the sequential method. If a level has an odd 
number of items, the leftover item is simply moved to the next level 
without being modified.

\paragraph{Atomicity of the reduce worker.}
A key implementation detail is that each pairwise merge executes as a single atomic unit inside one worker
process (Listing~\ref{lst:merge-worker}). This is what allows the merged node to be fed
back into the tree as an ordinary node at the next level, but it also means these four
sub-phases cannot be individually timed from the orchestrating process; only the aggregate
duration of the whole worker call is externally observable.

\begin{lstlisting}[language=Python, basicstyle=\small\ttfamily, breaklines=true, breakatwhitespace=true, frame=single, caption={Reduce phase worker: an atomic merge unit}, label={lst:merge-worker}]
def _merge_pair_worker(args):
    cv2.setNumThreads(1)
    (img_a, kp_a_ser, des_a), (img_b, kp_b_ser, des_b) = args

    kp_a = _deserialize_kp(kp_a_ser)
    kp_b = _deserialize_kp(kp_b_ser)

    matches = match_features(des_a, des_b)
    if len(matches) < 4:
        return (img_a, kp_a_ser, des_a)  # degrade gracefully

    H = estimate_homography(kp_a, kp_b, matches)
    if H is None:
        return (img_a, kp_a_ser, des_a)

    try:
        merged = warp_and_blend(img_a, img_b, H)
    except ValueError as e:
        # Degenerate homography detected (oversized canvas):
        # keep the left node unchanged rather than crash the worker.
        return (img_a, kp_a_ser, des_a)

    sift = cv2.SIFT_create(nfeatures=8000)
    gray = cv2.cvtColor(merged, cv2.COLOR_BGR2GRAY)
    kp_m, des_m = sift.detectAndCompute(gray, None)
    kp_m_ser = [(p.pt, p.size, p.angle, p.response,
                 p.octave, p.class_id) for p in kp_m]
    return (merged, kp_m_ser, des_m)
\end{lstlisting}

\paragraph{Correctness caveat.}
Unlike the linear-fold pipelines, this variant's output is \emph{not} expected to be
bit-identical to the sequential baseline. Composing homographies in tree order instead of
left-to-right changes which images are warped relative to which, and perspective
interpolation is inherently order-dependent; the resulting panorama can also differ in
overall shape when the reduce worker's fallback discards a pair. The panorama produced is
expected to remain geometrically plausible, but a pixel-level difference against the sequential
baseline is not a meaningful correctness signal for this pipeline.

\paragraph{A structural trade-off.}
The tree-merge structure introduces a subtler risk that the linear fold does not share.
In the linear fold, every new image is matched against the \emph{entire panorama
accumulated so far}, giving RANSAC a large, geometrically diverse pool of keypoints to draw
from. In the MapReduce tree, by contrast, each pair is fixed \emph{by position} in the
index sequence, not by verified spatial adjacency: if two images assigned to the same
tree node do not in fact overlap in the scene, no fallback within the algorithm itself
can discover an alternative pairing. Empirically, this trade-off manifests as a
substantially higher variance in the speedup achieved across image windows compared to
the other three parallel pipelines: while the linear-fold variants oscillate in a narrow
band, the MapReduce variant ranges from speedups exceeding $2\times$ down to occasional
slowdowns below $1\times$ when a fixed pairing turns out to be geometrically unfavorable.

\subsection{Joblib Implementation}

This fifth variant replaces the manual \texttt{ProcessPoolExecutor} orchestration used in
\texttt{parallel.py} with the high-level \texttt{joblib.Parallel} interface, retaining the
same overall architecture while delegating
process-pool lifecycle management to Joblib's \texttt{loky} backend.

\paragraph{Motivation.}
Every parallel and MapReduce variant discussed so far pays an inter-process communication
(IPC) cost whenever a NumPy image array crosses a process boundary: \texttt{multiprocessing}
pickles the full array and copies it into the receiving worker's address space. For a
half-resolution frame ($\approx 1.5$\,MB), this cost is paid once per image on every call to
the extraction phase, including every warm-up and timed repetition of the benchmark. Rather
than manage a \texttt{multiprocessing.shared\_memory} block by hand (the approach taken by the
shared-memory variant discussed in Section~\ref{sec:shared-memory}), this implementation
delegates the problem to Joblib's built-in automatic memmapping.

\paragraph{Automatic memmapping.}
When an array argument passed through \texttt{Parallel(...)(delayed(fn)(arg) for arg in
\ldots)} exceeds the \texttt{max\_nbytes} threshold, Joblib transparently writes it to a
temporary memory-mapped file on disk and passes only the file handle and shape/dtype metadata
to each worker, which maps it back into memory read-only with no additional copy. The 
threshold used here, \texttt{max\_nbytes="1M"}, is deliberately set 
well below the size of a single image frame (which is roughly 
6-8\,MB). This ensures that every large image automatically triggers 
memmapping, while small data payloads remain standard in-memory 
objects, avoiding unnecessary disk operations for tiny variables.

\begin{lstlisting}[language=Python, basicstyle=\small\ttfamily, breaklines=true, frame=single, breakatwhitespace=true, caption={Joblib-based parallel Map phase}, label={lst:joblib-map}]
_MEMMAP_THRESHOLD = "1M"
_BACKEND = "loky"

def extract_features_joblib(images):
    results = Parallel(
        n_jobs=NUM_CORES,
        backend=_BACKEND,
        max_nbytes=_MEMMAP_THRESHOLD,
        prefer="processes",
    )(delayed(_extract_worker)(img) for img in images)
    # ... deserialize keypoints, return (kp_list, des_list, extraction_time)
\end{lstlisting}

\paragraph{Backend and worker-pool reuse.} The \texttt{loky} backend 
is explicitly chosen using \texttt{backend="loky"} and 
\texttt{prefer="processes"}. This prevents Joblib from silently 
switching to a threading backend, ensuring that the CPU-heavy SIFT 
extraction uses separate processes and truly bypasses the GIL. 
Furthermore, \texttt{loky} automatically keeps its worker processes 
alive in the background across different calls. When 
\texttt{extract\_features\_joblib} is called repeatedly for different 
image windows, it reuses the same underlying processes instead of 
creating new ones from scratch every time. This avoids the high cost 
of spawning operating system processes, achieving the same efficiency 
as managing a long-lived \texttt{ProcessPoolExecutor} manually, but 
without the complex lifecycle code.

\subsection{Shared Memory Optimization}
\label{sec:shared-memory}
Inter-process communication (IPC) serialization is a major performance 
bottleneck in Python's multiprocessing model. By default, when a 
NumPy image array crosses a process boundary, the parent process 
pickles it and the receiving worker reconstructs it, copying the 
entire array for every transfer. For a full-resolution frame 
($\approx 4896 \times 3672 \times 3$ bytes, $\approx 54$\,MB), this 
copy cost is paid for every image during the MAP phase. In the 
MapReduce variant, this penalty occurs for every pair at every 
REDUCE level, where the transferred payload is an entire partial 
panorama instead of a single frame.

\paragraph{Strategy.}
This variant sidesteps the problem using \texttt{multiprocessing.shared\_memory}, which
allocates a contiguous block of raw bytes directly in the operating system's shared memory
space rather than pickling data through the multiprocessing IPC channel. After
\texttt{load\_images\_parallel()} produces the window's image list, each array is copied
\emph{once} into its own \texttt{SharedMemory} block; workers then receive only a
lightweight descriptor tuple (\texttt{shm\_name, shape, dtype}) instead of the
array itself, and reconstruct a NumPy view over the shared buffer with
\texttt{np.ndarray(shape, dtype=dtype, buffer=shm.buf)}.

\begin{lstlisting}[language=Python, basicstyle=\small\ttfamily, breaklines=true, frame=single, breakatwhitespace=true, caption={Allocating shared memory once for the whole window}, label={lst:shm-alloc}]
def _alloc_shm_for_images(images):
    blocks, descriptors = [], []
    for img in images:
        shm = SharedMemory(create=True, size=img.nbytes)
        buf_arr = np.ndarray(img.shape, dtype=img.dtype, buffer=shm.buf)
        np.copyto(buf_arr, img)          # one-time copy into shared RAM
        blocks.append(shm)
        descriptors.append((shm.name, img.shape, img.dtype.str))
    return blocks, descriptors
\end{lstlisting}

\paragraph{What is actually eliminated.}
It is important to state precisely which direction of communication this optimizes. Only
the \emph{outbound} transferf (from the main process to each worker) is avoided: a
worker no longer receives a pickled copy of the full image, only a small tuple of metadata.
The \emph{inbound} direction is unchanged: extracted keypoints and descriptors are still
returned to the main process through the ordinary \texttt{ProcessPoolExecutor} return
channel and pickled as usual, since they are orders of magnitude smaller than a raw frame
and not worth placing in shared memory themselves.

\paragraph{Not a genuine zero-copy read, by design.} A closer look at 
the worker function reveals that the read phase is not a strict 
zero-copy operation. Once connected, the worker explicitly copies the 
shared array into its local memory and immediately releases its 
connection to the shared block. This is a deliberate safety measure, 
not an oversight. If a worker kept the shared block open during the 
entire (and potentially slow) SIFT computation, the main process might 
try to clean up the memory (\texttt{shm.unlink()}) while the worker is 
still attached, causing the system to hang or crash. Making a quick 
local copy and releasing the connection keeps the shared memory's 
lifecycle short and safe, at the acceptable cost of one local memory 
copy per worker.


\section{Experimental Setup}
\label{sec:experimental-setup}

To ensure a statistically sound and reproducible comparison across the six stitching
pipelines described in the previous sections, a dedicated benchmarking framework
(\texttt{benchmark.py}) was developed. Rather than timing each pipeline in isolation with
ad-hoc scripts, every pipeline is wrapped in a uniform interface and driven through the same
sliding-window protocol, the same warm-up and repetition discipline, and the same statistical
post-processing, so that any observed difference in wall-clock time can be attributed to the
pipeline's architecture rather than to inconsistencies in how it was measured.

\subsection{Hardware and Software Environment}
All benchmarks were executed on a Windows machine equipped with an 
\textbf{Intel Core i7-8565U CPU @ 1.80\,GHz} (4 physical cores, 
8 logical threads) and \textbf{16\,GB} of RAM, running the 
CPython 3.14 interpreter. 

To maximize resource utilization, the framework is configured with 
the \texttt{NUM\_CORES} constant set to match the logical thread 
count of the evaluation hardware. This parameter strictly 
governs the worker count of every \texttt{ProcessPoolExecutor} and 
\texttt{ThreadPoolExecutor} instantiated throughout the experiments, 
as well as the number of horizontal tiles generated during the 
data-parallel blending phase.

\subsection{Pipeline Abstraction}
Each of the six variants (sequential, parallel, producer-consumer, MapReduce, Joblib, and
shared-memory) is described by a lightweight \texttt{PipelineSpec} record (Listing
\ref{lst:pipelinespec}): a display name, a \texttt{run(images, resources, seed)} callable
returning a dictionary of per-phase timings plus the resulting panorama, and two boolean
flags declaring whether the pipeline requires a process pool, a thread pool, or both. Before
benchmarking a given window, the framework inspects every selected \texttt{PipelineSpec} and
opens the \emph{union} of executors required by any of them exactly once, reusing the same
pools across the warm-up passes and all timed repetitions of every pipeline in that window.
This design avoids conflating executor spawn overhead with the architectural cost the benchmark is actually meant to
isolate.

\begin{lstlisting}[language=Python, basicstyle=\small\ttfamily, breaklines=true, frame=single, breakatwhitespace=true, caption={The pipeline abstraction driving the benchmark}, label={lst:pipelinespec}]
@dataclass
class PipelineSpec:
    name: str
    run: Callable[..., dict]
    needs_process_pool: bool = True
    needs_thread_pool: bool = False

pipelines_to_run = [
    SEQUENTIAL_SPEC,        # baseline
    PARALLEL_SPEC,
    PRODUCER_CONSUMER_SPEC,
    MAPREDUCE_SPEC,
    JOBLIB_SPEC,
    SHM_SPEC,
]
\end{lstlisting}

By convention, the first entry in the pipeline list is always the \emph{baseline}
(\texttt{sequential}); every subsequent entry is a \emph{candidate} reported and exported
against it, mirroring the sequential-vs-parallel comparison strategy used throughout this
study.

\subsection{Sliding-Window Protocol}
Because datasets can contain too many images to stitch into a single 
panorama, the benchmark operates on non-overlapping 
sliding windows of \texttt{WINDOW\_SIZE = 4} images. 

Images for each window are loaded from disk exactly once and kept in 
memory. They are reused unchanged across every warm-up pass and timed 
repetition of all six pipelines. Since no pipeline mutates the input 
images in place—each allocating a new panorama via 
\texttt{warp\_and\_blend}, this approach carries no correctness risk and 
completely eliminates disk I/O variance from the timed measurements. 

This design deliberately departs from a cold-cache protocol, where 
inputs would be reloaded before every single run. With six pipelines 
and multiple windows, avoiding repeated disk access drastically reduces 
the total benchmark duration.

\subsection{Warm-up, Repetitions, and Thermal Management}
For every pipeline, two untimed warm-up passes are executed first. 
This allows the underlying system and worker pools to reach a steady 
state before any measurements begin. After the warm-up, the benchmark 
records \texttt{N\_RUNS = 3} independent repetitions. 

To ensure consistent results, the framework forces an explicit 
garbage collection (\texttt{gc.collect()}) and pauses briefly between 
runs and between different pipelines. This approach clears leftover 
memory overhead and provides a short cooldown window for the CPU. 
By doing so, it prevents random garbage-collection pauses and hardware 
thermal throttling from skewing the final performance metrics.

\subsubsection{Statistical Methodology}
For every measured phase, the mean, standard deviation, and 95\% confidence interval are
computed using Student's \emph{t}-distribution, consistent with the
small-sample regime imposed by \texttt{N\_RUNS = 3} (2 degrees of freedom, two-tailed critical
value $t = 4.303$ at $\alpha = 0.05$). Speedup between the baseline and a candidate is then
computed as the ratio of their means, with its own confidence interval propagated via the
delta method, combining each side's coefficient of variation:
\[
\frac{\sigma_S}{S} \approx \sqrt{\left(\frac{\sigma_{\text{base}}}{\mu_{\text{base}}}\right)^2
+ \left(\frac{\sigma_{\text{cand}}}{\mu_{\text{cand}}}\right)^2}.
\]
Because a pipeline may complete fewer than \texttt{N\_RUNS} repetitions in a given window (see
Section \ref{subsec:fault_tolerance}), the baseline and candidate samples are permitted to differ in size;
each side's own sample size is used for its variance term, while the more conservative
(smaller) of the two sample sizes is used to look up the critical $t$-value, avoiding an
overly narrow interval when one side has fewer observations than the other.


\subsection{Handling Non-Measurable Phases}
Not every pipeline can report a clean, standalone duration for every phase of the stitching
pipeline. In the producer-consumer variant, extraction is deliberately
overlapped with the rest of the chain and has no isolated duration; in the MapReduce variant,
matching, homography, warping, and re-extraction all execute atomically inside a single
worker call per tree node and can't be timed separately from the orchestrating process. 
To handle this, such unmeasurable phases are directly recorded as \texttt{None}. 
This ensures they are explicitly excluded from all aggregate statistical 
computations, and they are neatly rendered as \texttt{"n/a"} in both the 
console reports and the exported CSV files.

\subsection{Correctness Validation}
After the timed runs for a window complete, the last panorama produced by 
each candidate is compared against the baseline on a pixel-by-pixel basis. 
The absolute per-pixel difference is computed using \texttt{int32} arithmetic 
to avoid \texttt{uint8} underflow. The framework evaluates the maximum and 
mean differences, the percentage of identical pixels, the per-channel BGR 
breakdown and the Peak Signal-to-Noise Ratio. 

A qualitative verdict is then assigned: \emph{bit-identical} when the 
maximum difference is zero; \emph{equivalent} when the difference is at most 
one least-significant bit due to floating-point rounding variations; or 
\emph{visually identical} when the Peak Signal-to-Noise Ratio exceeds 40\,dB. 
Outcomes failing these criteria are flagged for manual inspection. Shape 
mismatches, which are expected for the MapReduce variant, are reported as 
distinct outcomes to avoid meaningless pixel comparisons.

\subsection{Robustness and Fault Tolerance}
\label{subsec:fault_tolerance}
The benchmarking framework implements two safeguards to prevent isolated errors from invalidating the entire execution. 

First, each call to the warping and blending routine is wrapped in a \texttt{try/except} block. 
This block catches a \texttt{ValueError} which is triggered if a degenerate homography attempts to create an oversized canvas. 
The framework defines this as any canvas exceeding four times the combined input area. 
When this occurs, the problematic image is skipped to prevent the system from exhausting its memory. 

Second, all warm-up passes and timed repetitions are individually protected by exception handlers. 
If a pipeline fails during its warm-up or across all its timed runs within a specific window, it is simply excluded from that window's results. 
If the baseline pipeline fails, the framework immediately aborts the remainder of that window since comparisons are no longer possible. 
However, all previously recorded results for other windows and pipelines remain safely preserved.

\subsection{Output Artifacts}
For each window, the framework exports two CSV files under \texttt{results/}:
\texttt{benchmark\_results.csv}, containing per-phase mean, standard deviation, margin, and
speedup with confidence bounds for every (baseline, candidate) pair and
\texttt{correctness\_results.csv}, containing the pixel-comparison metrics described above.
The last panorama produced by every pipeline in every window is additionally saved to
\texttt{data/output/benchmark/} for visual inspection, independent of the numerical results.

{\small
\bibliographystyle{ieee.bst}
\bibliography{egbib.bib}
}



\end{document}
